\documentclass{article}
\usepackage[utf8]{inputenc}
\usepackage[absolute]{textpos}
\usepackage[default]{raleway}
\usepackage{titlesec, comment, tabularx, makecell, listings, array, setspace, geometry, graphicx, xcolor, xparse, fancyvrb, relsize, fancyhdr, booktabs, hyperref}
\usepackage{colortbl}
\usepackage{makecell}
%\geometry{a4paper, left=2cm, right=2cm, top=2cm, bottom=2.5cm}
\renewcommand{\headrulewidth}{0pt}

% Definisci uno stile per i comandi git
\definecolor{light-gray}{gray}{0.92}

\lstdefinestyle{code}{
    frame=single,
    framesep=1mm,
    rulecolor=\color{light-gray},
    backgroundcolor=\color{light-gray},
    basicstyle=\ttfamily,
}

% ----------------------------- Definizione tabella ----------------------------

\newcolumntype{C}[1]{>{\centering\arraybackslash}m{#1}}

%\setcellgapes{2ex} % Imposta l'altezza dell'header (2ex)


% ------------------------------Metadati indice --------------------------------
\title{\textbf{\fontsize{28}{6}\selectfont Indice}}
\author{\fontsize{14}{6}\selectfont ByteOps}
\date{Ottobre 20, 2023}


% -----------------------------Creazione footer --------------------------------

\pagestyle{fancy}
\fancyhf{}
\renewcommand{\footrulewidth}{0.4pt}
\lfoot{
    \parbox[c]{2cm}{\includegraphics[width=2cm]{../../Images/logo.png}}
    \textcolor[RGB]{120, 120, 120}{$\cdot$ Verbale Esterno}
}
\rfoot{\thepage}

% --------------------------Modifica formato hyperlinks ------------------------

\hypersetup{
    colorlinks=true,
    linkcolor=black,
    filecolor=black,      
    pdftitle={Verbale 20-10-2023},
    pdfpagemode=FullScreen,
}

% ------------------------------- Valore sotto-paragrafi indice --------------------------------------

\setcounter{secnumdepth}{4}
\setcounter{tocdepth}{4}

\titleformat{\section}
{\normalfont\huge\bfseries}{\thesection}{0.2cm}{}
\titlespacing*{\paragraph}{0pt}{0.5cm}{0.1cm}

\titleformat{\subsection}
{\normalfont\Large\bfseries}{\thesubsection}{0.2cm}{}
\titlespacing*{\paragraph}{0pt}{0.5cm}{0.1cm}

\titleformat{\subsubsection}
{\normalfont\large\bfseries}{\thesubsubsection}{0.2cm}{}
\titlespacing*{\paragraph}{0pt}{0.5cm}{0.1cm}

\titleformat{\paragraph}
{\normalfont\normalsize\bfseries}{\theparagraph}{0.2cm}{}
\titlespacing*{\paragraph}{0pt}{0.5cm}{0.1cm}

% ------------------------------- Front Page ---------------------------------------

\begin{document}

% --------------------------Aggiunta firma finale ------------------------
\begin{textblock*}{\textwidth}(0.85\textwidth, 1.16\textheight)
    Il responsabile: Davide Diotto
\end{textblock*}
% ------------------------------------------------------------------------

\pagestyle{fancy}
\begin{center}
\includegraphics[width = 0.7\textwidth]{../../Images/logo.png} \\
\vspace{0.2cm}
\textcolor[RGB]{60, 60, 60}{\textit{ByteOps.swe@gmail.com}} \\
\vspace{1cm}
\fontsize{16}{6}\selectfont Verbale Esterno $\cdot$ Data: 20/10/2023 \\
\vspace{0.5cm}
\end{center}

\section*{Informazioni documento}
\def\arraystretch{1.2} \begin{tabular}{>{\raggedleft\arraybackslash}p{0.2\textwidth}|>{\raggedright\arraybackslash}p {0.6\textwidth}c}
\hline
\addlinespace
\textbf{Luogo} & Goole Meet \vspace{10pt} \\
\textbf{Orario} & 11:30 - 12:00 \vspace{10pt} \\
\textbf{Redattori} & A. Barutta \\ & R. Smanio \\ & N. Preto 
\vspace{10pt} \\ 
\textbf{Verificatori} & E. Hysa \\ & L. Skenderi \\ & D. Diotto 
\vspace{10pt} \\ 
\textbf{Amministratore} & F. Pozza 
\vspace{10pt} \\
\textbf{Destinatari} & T. Vardanega \\ & R. Cardin \vspace{10pt} \\
\textbf{Partecipanti} & A. Barutta \\ & E. Hysa \\ & R. Smanio \\ & D. Diotto \\ & F. Pozza \\ & L. Skenderi \\ & N. Preto \\ & A. Dorigo, D. Zorzi, F. Pallaro \vspace{10pt} \\
\end{tabular}
\pagebreak 


% ------------------------- Changelog ----------------------------
\section*{Registro delle modifiche}
\begin{tabular}{|C{2.5cm}|C{2.5cm}|C{2.5cm}|C{2.5cm}|C{2.5cm}|} \hline
\textbf{Versione} & \textbf{Data} & \textbf{Autore} & \textbf{Verificatore} & \textbf{Dettaglio} \\
\hline \hline
0.0.2 & 07/11/2023 & A. Barutta & \makecell{E. Hysa\\L. Skenderi} & Nuovo template verbali \\ \hline
0.0.1 & 20/10/2023 & \makecell{A. Barutta\\R. Smanio} & D. Diotto & Redazione documento\\ \hline
\end{tabular} % \pagebreak

% ------------------------- Generazione automatica indice ----------------------
\setstretch{1.5}
\maketitle
\thispagestyle{fancy}
\tableofcontents 
\setstretch{1.2} 
\pagebreak

% ------------------------ INIZIO DOCUMENTO ----------------------
\flushleft

\section{Revisione del periodo precedente} 
    La presente riunione rappresenta il nostro primo incontro ufficiale con l'azienda Sync Lab, \textit{committente}\textsubscript{\textit{G}} del capitolato C6. Questa fase iniziale del potenziale rapporto di collaborazione è stata concepita con l'intento di instaurare una conoscenza reciproca e risolvere eventuali incertezze. Pertanto, non è stato possibile effettuare revisioni dettagliate delle \textit{attività}\textsubscript{\textit{G}} svolte fino a quel momento.  

\section{Ordine del giorno} 
    \subsection{Modalità di comunicazione proponente/fornitore}
        Verrà fornito da Sync Lab un invito per entrare all’interno di un server \textit{Discord}\textsubscript{\textit{G}}.
        Per le videoconferenze si utilizzerà Google Meet. 
        Verrà pubblicato un \textit{calendario}\textsubscript{\textit{G}} condiviso tramite Google Drive per essere a conoscenza in anticipo dei vari incontri con l’azienda \textit{proponente}\textsubscript{\textit{G}}. Un’alternativa proposta è utilizzare anche Google Calendar per essere in grado di monitorare meglio i vari gruppi. 

    \subsection{SAL periodici}
        Saranno organizzati, con probabile, ma non ancora definita cadenza settimanale, degli incontri periodici per monitorare e supportare il corretto avanzamento del progetto.
        Gli obiettivi tra un \textit{SAL}\textsubscript{\textit{G}} e il successivo saranno ben definiti e saranno determinati con il supporto di Sync Lab. 

    \subsection{Approfondimento sulla provenienza dei dati generati}
        Il dubbio nasce dalla richiesta di fornire dati al \textit{sistema}\textsubscript{\textit{G}} non solo da sensori da noi simulati con \textit{script}\textsubscript{\textit{G}}, ma anche attraverso altre fonti esterne non specificate.
        I dati, oltre ai sensori, possono provenire da diverse fonti come ad esempio dispositivi indossabili e auto. 
        Inoltre è stato precisato che il focus non è sul simulatore poiché serve solo per testare l’efficacia del prodotto.
        Sarà possibile discutere su come generare la fonte dei dati. L’azienda è disponibile a fornirci un set di dati da cui prendere spunto.
        Non viene consigliato l’uso di \textit{API}\textsubscript{\textit{G}} per il recupero di dati forniti da sensori realmente in uso perché si necessita di avere un elevato numero di dati, cosa che l’\textit{API}\textsubscript{\textit{G}} non garantisce. \subsection{Pulizia di eventuali dati errati} 
        Nella realtà i dati raccolti dai sensori vengono filtrati. Non è un requisito obbligatorio, ma molto gradito, sul quale ci potrebbe essere una futura contrattazione. \subsection{Chiarimenti a riguardo: “Strumenti proposti nel capitolato”} 
        Verranno organizzati dei meeting su Google Meet (o anche in presenza nella loro sede) per spiegare i concetti fondamentali su cui si basano le tecnologie suggerite. 
    
    \subsection{Cosa si intende per test end to end}
        Si intende che l’applicazione deve essere testata nella sua interezza,”dall’inizio alla fine”, simulando l’interazione da parte del cliente/utente finale. Ciò permette di verificare che l’applicazione si comporti come ci si aspetta. 

    \subsection{Chiarimenti a riguardo: “script python generano stream di messaggi, i tracciati devono essere proposti dal fornitore e validati dal proponente”}
        Con “tracciati” si intende un file in cui vengono inseriti i dati raccolti dai sensori (formato JSON consigliato). Viene consigliato \textit{Python}\textsubscript{\textit{G}} perché considerato più adatto e semplice per la realizzazione degli \textit{script}\textsubscript{\textit{G}}. 
    
    \subsection{Way of working proposto dal proponente}
        L’azienda fissa un \textit{calendario}\textsubscript{\textit{G}} in anticipo con dei \textit{SAL}\textsubscript{\textit{G}} e degli obiettivi da raggiungere (di solito su base settimanale).
        Ogni \textit{SAL}\textsubscript{\textit{G}} permette di verificare lo stato di avanzamento del progetto. 
        È molto importante non “buttarsi” subito nello sviluppo, ma è necessario prestare particolare attenzione all’organizzazione e alla pianificazione delle \textit{attività}\textsubscript{\textit{G}} per perseguire efficienza ed efficacia. 

\section{Attività da svolgere}
    Non sono state definite nuove \textit{attività}\textsubscript{\textit{G}}. 

\end{document}
